\section{复算、独立检查与覆盖分析}

\subsection{结果可追溯性}

运行包把输入配置、源程序、完整 JSONL 观测、Metric、图表数据、绘图脚本、显示配置和检查结果放在同一
版本化目录中。发布阶段首先核对正式运行标识与清单 SHA-256，再逐行比较图表 CSV 和完整观测；只有二者
数值一致时才重绘图。表~\ref{tab:trace} 给出论文对象到支撑材料的映射。

\small
\begin{longtable}{p{0.18\linewidth}p{0.30\linewidth}p{0.42\linewidth}}
  \caption{论文结论到运行证据的追溯关系}\label{tab:trace}\\
  \toprule
  论文对象 & 直接来源 & 发布检查\\
  \midrule
  \endfirsthead
  \toprule
  论文对象 & 直接来源 & 发布检查\\
  \midrule
  \endhead
  两个方法均值 & \path{run_manifest.json} 的 metrics & 定义、单位、聚合、分母与数值共同生成 TeX 宏\\
  实例级数值 & \path{observations.jsonl} & 图表 CSV 的每一行必须与完整观测相等\\
  计划覆盖 & manifest 的 instances & 分别统计 completed、failed 和 planned，不把失败当作零\\
  已知最优上界 & \path{checks/known-optimum-bound.json} & 由独立于 Metric 重算的检查实现验证\\
  图形表达 & figure data / script / display config & 只改变尺寸、颜色、marker、线型、hatch 和布局\\
  引用与匿名 & TeX、BibTeX、最终 PDF 文本层 & 检查 citekey、禁用身份词和第一页内容\\
  \bottomrule
\end{longtable}
\normalsize

\subsection{确定性复算}

由两个完整观测直接得到
\begin{align}
  \bar V_G &= \frac{12+21}{2}=\BaselineMeanValue,\\
  \bar V_E &= \frac{18+21}{2}=\AlternativeMeanValue.
\end{align}
这与清单中的 \texttt{metric\_recompute} 检查一致。另一个独立检查验证穷举结果没有超过配置中的已知最优值，
其职责与均值重算不同：前者检查约束/界，后者只聚合已保存观测。

\subsection{部分结果的解释}

第三个实例在真实子进程尝试后失败，因此状态是 \texttt{failed}，不是 \texttt{not\_run}。
本演示不删除该实例，也不把失败值补成 $0$。覆盖率可写为
\begin{equation}
  \rho_c=\frac{n_c}{n_p}=\frac{\RunCompleted}{\RunPlanned}=\frac{2}{3},
  \qquad
  \rho_f=\frac{n_f}{n_p}=\frac{\RunFailed}{\RunPlanned}=\frac{1}{3}.
\end{equation}
由于只有两个确定性完整观测，无法从运行包定义标准误、置信区间或随机种子间波动。因此图中没有误差棒，
均值差 $3$ value 也只作这两个完整实例的描述，不作为总体显著性判断。
